\subsection{\label{sec.sf.m}$N$-partitions and monomial symmetric polynomials}

Recall that an \emph{(integer) partition} means a weakly decreasing finite
tuple of positive integers -- such as $\left(  5,3,3,2,1\right)  $.

Let us define a variant of this notion:

\begin{definition}
\label{def.sf.Npar}An $N$\emph{-partition} will mean a weakly decreasing
$N$-tuple of nonnegative integers. In other words, an $N$-partition means an
$N$-tuple $\left(  \lambda_{1},\lambda_{2},\ldots,\lambda_{N}\right)
\in\mathbb{N}^{N}$ with $\lambda_{1}\geq\lambda_{2}\geq\cdots\geq\lambda_{N}$.
\end{definition}

For example, $\left(  5,3,3,2,1,0,0\right)  $ is a $7$-partition.

Per se, an $N$-partition can contain zeroes and thus is not always a
partition. However, the $N$-partitions are \textquotedblleft more or less the
same\textquotedblright\ as the partitions of length $\leq N$. Indeed:

\begin{proposition}
\label{prop.sf.Npar-as-par}There is a bijection%
\begin{align*}
\left\{  \text{partitions of length }\leq N\right\}   &  \rightarrow\left\{
N\text{-partitions}\right\}  ,\\
\left(  \lambda_{1},\lambda_{2},\ldots,\lambda_{\ell}\right)   &
\mapsto\left(  \lambda_{1},\lambda_{2},\ldots,\lambda_{\ell}%
,\underbrace{0,0,\ldots,0}_{N-\ell\text{ zeroes}}\right)  .
\end{align*}

\end{proposition}

\begin{proof}
Straightforward. (We essentially did this back in our proof of Proposition
\ref{prop.pars.qbinom.alt-defs} \textbf{(a)}, although we used the letter $k$
instead of $N$ back then.)
\end{proof}

The $N$-partitions turn out to be closely connected to the ring $\mathcal{S}$.
Indeed, we will soon see various bases of the $K$-module $\mathcal{S}$, all of
which are indexed by the $N$-partitions. We shall construct the simplest one
in a moment. First, we define some auxiliary notations:

\begin{definition}
\label{def.sf.sort}Let $a=\left(  a_{1},a_{2},\ldots,a_{N}\right)
\in\mathbb{N}^{N}$. Then: \medskip

\textbf{(a)} We let $x^{a}$ denote the monomial $x_{1}^{a_{1}}x_{2}^{a_{2}%
}\cdots x_{N}^{a_{N}}$. \medskip

\textbf{(b)} We let $\operatorname*{sort}a$ mean the $N$-partition obtained
from $a$ by sorting the entries of $a$ in weakly decreasing order.
\end{definition}

For example, if $N=5$, then $x^{\left(  1,5,0,4,4\right)  }=x_{1}^{1}x_{2}%
^{5}x_{3}^{0}x_{4}^{4}x_{5}^{4}=x_{1}x_{2}^{5}x_{4}^{4}x_{5}^{4}$ and
\newline$\operatorname*{sort}\left(  1,5,0,4,4\right)  =\left(
5,4,4,1,0\right)  $.

\begin{definition}
\label{def.sf.m}Let $\lambda$ be any $N$-partition. Then, we define a
symmetric polynomial $m_{\lambda}\in\mathcal{S}$ by%
\[
m_{\lambda}:=\sum_{\substack{a\in\mathbb{N}^{N};\\\operatorname*{sort}%
a=\lambda}}x^{a}.
\]
This is called the \emph{monomial symmetric polynomial corresponding to
}$\lambda$.
\end{definition}

\begin{example}
Let $N=3$. Then,%
\begin{align*}
m_{\left(  2,1,0\right)  }  &  =\sum_{\substack{a\in\mathbb{N}^{3}%
;\\\operatorname*{sort}a=\left(  2,1,0\right)  }}x^{a}=x^{\left(
2,1,0\right)  }+x^{\left(  2,0,1\right)  }+x^{\left(  1,2,0\right)
}+x^{\left(  1,0,2\right)  }+x^{\left(  0,2,1\right)  }+x^{\left(
0,1,2\right)  }\\
&  =x_{1}^{2}x_{2}+x_{1}^{2}x_{3}+x_{1}x_{2}^{2}+x_{1}x_{3}^{2}+x_{2}^{2}%
x_{3}+x_{2}x_{3}^{2}%
\end{align*}
and%
\begin{align*}
m_{\left(  3,2,1\right)  }  &  =\sum_{\substack{a\in\mathbb{N}^{3}%
;\\\operatorname*{sort}a=\left(  3,2,1\right)  }}x^{a}=x^{\left(
3,2,1\right)  }+x^{\left(  3,1,2\right)  }+x^{\left(  2,3,1\right)
}+x^{\left(  2,1,3\right)  }+x^{\left(  1,3,2\right)  }+x^{\left(
1,2,3\right)  }\\
&  =x_{1}^{3}x_{2}^{2}x_{3}+x_{1}^{3}x_{2}x_{3}^{2}+x_{1}^{2}x_{2}^{3}%
x_{3}+x_{1}^{2}x_{2}x_{3}^{3}+x_{1}x_{2}^{3}x_{3}^{2}+x_{1}x_{2}^{2}x_{3}%
^{3}\\
&  =x_{1}^{3}x_{2}^{2}x_{3}+\left(  \text{all other }5\text{ permutations of
this monomial}\right)
\end{align*}
and%
\begin{align*}
m_{\left(  2,2,1\right)  }  &  =\sum_{\substack{a\in\mathbb{N}^{3}%
;\\\operatorname*{sort}a=\left(  2,2,1\right)  }}x^{a}=x^{\left(
2,2,1\right)  }+x^{\left(  2,1,2\right)  }+x^{\left(  1,2,2\right)  }\\
&  =x_{1}^{2}x_{2}^{2}x_{3}+x_{1}^{2}x_{2}x_{3}^{2}+x_{1}x_{2}^{2}x_{3}^{2}%
\end{align*}
and%
\[
m_{\left(  2,2,2\right)  }=\sum_{\substack{a\in\mathbb{N}^{3}%
;\\\operatorname*{sort}a=\left(  2,2,2\right)  }}x^{a}=x_{1}^{2}x_{2}^{2}%
x_{3}^{2}.
\]

\end{example}

Our symmetric polynomials $e_{n}$, $h_{n}$ and $p_{n}$ so far can be easily
expressed in terms of monomial symmetric polynomials:

\begin{proposition}
\label{prop.sf.ehp-through-m}\textbf{(a)} For each $n\in\left\{
0,1,\ldots,N\right\}  $, we have%
\[
e_{n}=m_{\left(  1,1,\ldots,1,0,0,\ldots,0\right)  },
\]
where $\left(  1,1,\ldots,1,0,0,\ldots,0\right)  $ is the $N$-tuple that
begins with $n$ many $1$'s and ends with $N-n$ many $0$'s. \medskip

\textbf{(b)} For each $n\in\mathbb{N}$, we have%
\[
h_{n}=\sum_{\substack{\lambda\text{ is an }N\text{-partition;}\\\left\vert
\lambda\right\vert =n}}m_{\lambda},
\]
where the \emph{size} $\left\vert \lambda\right\vert $ of an $N$-partition
$\lambda$ is defined to be the sum of its entries (i.e., if $\lambda=\left(
\lambda_{1},\lambda_{2},\ldots,\lambda_{N}\right)  $, then $\left\vert
\lambda\right\vert :=\lambda_{1}+\lambda_{2}+\cdots+\lambda_{N}$). \medskip

\textbf{(c)} Assume that $N>0$. For each $n\in\mathbb{N}$, we have%
\[
p_{n}=m_{\left(  n,0,0,\ldots,0\right)  },
\]
where $\left(  n,0,0,\ldots,0\right)  $ is the $N$-tuple that begins with an
$n$ and ends with $N-1$ zeroes.
\end{proposition}

\begin{proof}
Easy and LTTR.
\end{proof}

The monomial symmetric polynomials $m_{\lambda}$ are the \textquotedblleft
building blocks\textquotedblright\ of symmetric polynomials, in the same way
as the monomials are the \textquotedblleft building blocks\textquotedblright%
\ of polynomials. Here is a way to make this precise:

\begin{theorem}
\label{thm.sf.m-basis}\textbf{(a)} The family $\left(  m_{\lambda}\right)
_{\lambda\text{ is an }N\text{-partition}}$ is a basis of the $K$-module
$\mathcal{S}$. \medskip

\textbf{(b)} Each symmetric polynomial $f\in\mathcal{S}$ satisfies%
\[
f=\sum_{\substack{\lambda=\left(  \lambda_{1},\lambda_{2},\ldots,\lambda
_{N}\right)  \\\text{is an }N\text{-partition}}}\left(  \left[  x_{1}%
^{\lambda_{1}}x_{2}^{\lambda_{2}}\cdots x_{N}^{\lambda_{N}}\right]  f\right)
m_{\lambda}.
\]


\textbf{(c)} Let $n\in\mathbb{N}$. Let%
\[
\mathcal{S}_{n}:=\left\{  \text{homogeneous symmetric polynomials }%
f\in\mathcal{P}\text{ of degree }n\right\}
\]
(where we understand the zero polynomial $0\in\mathcal{P}$ to be homogeneous
of every degree). Then, $\mathcal{S}_{n}$ is a $K$-submodule of $\mathcal{S}$.
\medskip

\textbf{(d)} Define the \emph{size} of any $N$-partition $\lambda=\left(
\lambda_{1},\lambda_{2},\ldots,\lambda_{N}\right)  $ to be the number
$\lambda_{1}+\lambda_{2}+\cdots+\lambda_{N}\in\mathbb{N}$. Then, the family
$\left(  m_{\lambda}\right)  _{\lambda\text{ is an }N\text{-partition of size
}n}$ is a basis of the $K$-module $\mathcal{S}_{n}$.
\end{theorem}

\begin{example}
\label{exa.sf.m-basis.1}Let $N=3$, and let us rename the indeterminates
$x_{1},x_{2},x_{3}$ as $x,y,z$. The polynomial $\left(  x+y\right)  \left(
y+z\right)  \left(  z+x\right)  $ is symmetric, thus belongs to $\mathcal{S}$.
Expanding it, we find%
\begin{align*}
\left(  x+y\right)  \left(  y+z\right)  \left(  z+x\right)   &
=\underbrace{x^{2}y+x^{2}z+y^{2}x+y^{2}z+z^{2}x+z^{2}y}_{=m_{\left(
2,1,0\right)  }}+\,2\underbrace{xyz}_{=m_{\left(  1,1,1\right)  }}\\
&  =m_{\left(  2,1,0\right)  }+2m_{\left(  1,1,1\right)  }.
\end{align*}
Thus, we have written $\left(  x+y\right)  \left(  y+z\right)  \left(
z+x\right)  $ as a $K$-linear combination of $m_{\lambda}$'s for various
$N$-partitions $\lambda$. The same procedure (i.e., expanding, and then
collecting monomials that differ only in the order of their exponents, such as
the monomials $x^{2}y,x^{2}z,y^{2}x,y^{2}z,z^{2}x,z^{2}y$ in our example) can
be applied to any symmetric polynomial $f\in\mathcal{S}$, and always results
in a representation of $f$ as a $K$-linear combination of $m_{\lambda}$'s
(because the symmetry of $f$ ensures that monomials that differ only in the
order of their exponents appear in $f$ with equal coefficients).
\end{example}

The proof of Theorem \ref{thm.sf.m-basis} will rely on a simple proposition
that expresses how a permutation $\sigma\in S_{N}$ transforms the coefficients
of a polynomial $f\in\mathcal{P}$ (guess what: it permutes these coefficients):

\begin{proposition}
\label{prop.sf.sigma-pol-coeff}Let $\sigma\in S_{N}$ and $f\in\mathcal{P}$.
Then,%
\[
\left[  x_{1}^{a_{1}}x_{2}^{a_{2}}\cdots x_{N}^{a_{N}}\right]  \left(
\sigma\cdot f\right)  =\left[  x_{1}^{a_{\sigma\left(  1\right)  }}%
x_{2}^{a_{\sigma\left(  2\right)  }}\cdots x_{N}^{a_{\sigma\left(  N\right)
}}\right]  f
\]
for any $\left(  a_{1},a_{2},\ldots,a_{N}\right)  \in\mathbb{N}^{N}$.
\end{proposition}

Here, as in Section \ref{sec.gf.multivar}, we are using the notation $\left[
x_{1}^{a_{1}}x_{2}^{a_{2}}\cdots x_{N}^{a_{N}}\right]  g$ for the coefficient
of a monomial $x_{1}^{a_{1}}x_{2}^{a_{2}}\cdots x_{N}^{a_{N}}$ in a polynomial
$g\in\mathcal{P}$.

The proof of Proposition \ref{prop.sf.sigma-pol-coeff} is quite easy and can
be found in Section \ref{sec.details.sf.m}.

\begin{proof}
[Proof of Theorem \ref{thm.sf.m-basis} (sketched).]Here is a rough outline of
the proof; we leave the details to the reader. \medskip

\textbf{(a)} The method shown in Example \ref{exa.sf.m-basis.1} shows that
each $f\in\mathcal{S}$ is a $K$-linear combination of the family $\left(
m_{\lambda}\right)  _{\lambda\text{ is an }N\text{-partition}}$. Thus, this
family spans $\mathcal{S}$. It remains to show that this family is
$K$-linearly independent.

To show this, we observe the following: If you expand a linear combination
$\sum\limits_{\lambda\text{ is an }N\text{-partition}}a_{\lambda}m_{\lambda}$
(where $a_{\lambda}\in K$) as a sum of monomials, then none of the addends can
cancel (since no two $m_{\lambda}$'s have any monomial in common\footnote{and
since each $m_{\lambda}$ contains at least one monomial (trivial observation,
but should not be forgotten)}); thus, the linear combination cannot be $0$
unless all the $a_{\lambda}$'s are $0$. This proves the $K$-linear
independence of the family $\left(  m_{\lambda}\right)  _{\lambda\text{ is an
}N\text{-partition}}$. Thus, the proof of Theorem \ref{thm.sf.m-basis}
\textbf{(a)} is complete. \medskip

\textbf{(b)} This should be clear from Example \ref{exa.sf.m-basis.1} as well.
\medskip

\textbf{(c)} This is rather obvious: Any $K$-linear combination of homogeneous
polynomials of degree $n$ is again homogeneous of degree $n$. \medskip

\textbf{(d)} This follows by the same argument as part \textbf{(a)}, except
that we now need to observe that homogeneous polynomials of degree $n$ are
$K$-linear combinations of degree-$n$ monomials (rather than arbitrary monomials).
\end{proof}
